% Context from chapters/wavelets.tex; for mathematical reference, not compilation.
resolution Approximation Spaces}

A multiresolution approximation of $L^2(\RR)$ consists of nested closed subspaces $(V_j)_j$
\eql{\label{eq-nested-spaces}
	L^2(\RR) \supset \ldots \supset V_{j-1} \supset V_j \supset V_{j+1} \supset \ldots \supset \{0\}
}
related by dyadic scaling and invariant under translations on the corresponding dyadic grid:
\eq{
	f \in V_j \quad\Longleftrightarrow\quad f(\cdot/2) \in V_{j+1}
	\qandq
	f \in V_j \quad\Longleftrightarrow\quad \foralls n \in \ZZ, \: f(\cdot + n 2^j) \in V_{j}
}
Large $j$ corresponds to coarse approximation spaces, and $2^j$ is often called the ``scale''.

At the fine-scale limit on the left of~\eqref{eq-nested-spaces}, $\cup_j V_j$ is dense in $L^2(\RR)$. Equivalently, $P_{V_j}(f) \rightarrow f$ as $j \rightarrow -\infty$, where $P_V$ denotes orthogonal projection onto $V$:
\eq{
	P_V(f) = \uargmin{f' \in V} \norm{f-f'}.
}
The limit on the right of~\eqref{eq-nested-spaces} means that $\cap_j V_j = \{0\}$, or equivalently that $P_{V_j}(f) \rightarrow 0$ as $j \rightarrow +\infty$.

The first example consists of piecewise constant functions on dyadic intervals
\eql{\label{eq-haar-multires}
	V_j = \enscond{f \in L^2(\RR)}{\forall n, f \text{ is constant on } [2^j n, 2^j (n+1)[}, 
}

A second example is the space used for Shannon interpolation of bandlimited signals
\eql{\label{eq-shannon-multires}
	V_j = \enscond{f}{ \Supp(\hat f) \subset [-2^{-j}\pi,2^{-j} \pi] } 
}
whose elements are recovered exactly from the samples $f(2^j n)$. As for piecewise constant signals, the sampling grid has spacing $2^j$. Here, however, orthogonal projection is a bandlimiting operation:
\eq{
	P_{V_j}(f) = \Ff^{-1}( \hat f \odot 1_{[-2^{-j}\pi,2^{-j} \pi]} ). 
}

  
\paragraph{Scaling functions.}

We also require a scaling function $\phi \in L^2(\RR)$ such that
\eq{
	\{\phi(\cdot-n) \}_n  
	\text{ is a Hilbertian orthonormal basis of }
	V_0.
}
By the dilation property, this implies that 
\eq{
	\{ \phi_{j,n} \}_n \text{ is a Hilbertian orthonormal basis of } V_j
	\qwhereq
	\phi_{j,n} \eqdef \frac{1}{2^{j/2}} \phi\pa{ \frac{\cdot - 2^j n}{2^j} }.
}
The normalization ensures $\norm{\phi_{j,n}} = 1$.

Note that one then has
\eq{
	P_{V_j}(f) = \sum_n \dotp{f}{\phi_{j,n}} \phi_{j,n}. 
}
Figure~\ref{fig-multires-1d} illustrates the effects of translation and scaling.


\begin{figure}[tbp]
\centering
\BookPanel{.315}{1}{tikz/wavelets-scaling-1}\hfill
\BookPanel{.315}{2}{tikz/wavelets-scaling-2}\hfill
\BookPanel{.315}{3}{tikz/wavelets-scaling-3}
\caption{Translations and dilations of a scaling function generate the approximation spaces. The Shannon scaling function $\phi(t)=\sin(\pi t)/(\pi t)$ is illustrated, with horizontal coordinate $x/2^j$ and heights relative to $2^{-j/2}$.}
\label{fig-multires-1d}
\end{figure}


For the case of piecewise constant signals~\eqref{eq-haar-multires}, one can use 
\eq{
	\phi = 1_{[0,1[}
	\qandq 
	\phi_{j,n} = 2^{-j/2}1_{[2^j n, 2^j (n+1)[}.
}

For the case of Shannon multiresolution~\eqref{eq-shannon-multires}, one can use 
$\phi(t) = \sin(\pi t)/(\pi t)$ and one verifies
\eq{
	\dotp{f}{\phi_{j,n}}=2^{j/2}f(2^jn)\quad(f\in V_j)
	\qandq
	f=\sum_n2^{j/2}f(2^jn)\phi_{j,n}\quad(f\in V_j).
}

\myfigure{
\begin{minipage}{.9\linewidth}\centering
$\displaystyle\phi_{j,n}(x)=2^{-j/2}\mathbf1_{[2^jn,\,2^j(n+1))}(x)$\\[1em]
$\displaystyle\psi_{j,n}(x)=2^{-j/2}\left(\mathbf1_{[2^jn,\,2^j(n+1/2))}(x)-\mathbf1_{[2^j(n+1/2),\,2^j(n+1))}(x)\right)$
\end{minipage}
}{Haar scaling and wavelet functions. Both have amplitude $2^{-j/2}$ and unit $L^2$ norm.}{fig-haard-scaling-wav}

  
\paragraph{Spectral orthogonalization.}

In many applications, $V_0$ is generated by translates of a function that are not orthogonal: $V_0=\overline{\Span}\{\th(\cdot-n)\}_{n\in\ZZ}$. The next proposition orthogonalizes this family in the Fourier domain.
 
Figure~\ref{fig-spline-ortho} shows the orthogonal cardinal spline functions obtained by this construction.

\begin{prop}\label{prop-spectral-ortho}
	Let $\th \in L^2(\RR)$. Its translates $\{\th(\cdot-n)\}_{n \in \ZZ}$ are orthonormal if and only if
	\eq{
		A(\om) \eqdef \sum_{k} |\hat \th(\om-2\pi k)|^2 = 1
		\quad\text{for almost every }\om.
	}
	If there exist $0 < a \leq b <+\infty$ such that $a \leq A(\om) \leq b$ almost everywhere, then $\phi$ defined by
	\eq{
		\hat\phi(\om) = \frac{\hat\th(\om)}{\sqrt{A(\om)}}
	}
	is such that $\{\phi(\cdot-n)\}_{n \in \ZZ}$ is an orthonormal basis of $\overline{\Span}\{\th(\cdot-n)\}_{n\in\ZZ}$.
\end{prop}


\begin{figure}[tbp]
\centering
\BookDrawing{tikz/wavelets-spline-orthogonalization}
\caption{Orthogonalization of B-splines to define cardinal orthogonal spline functions.}
\label{fig-spline-ortho}
\end{figure}

\begin{proof}
	The family $\{\th(\cdot-n)\}_{n \in \ZZ}$ is orthonormal if and only if
	\eq{
		\dotp{\th}{\th(\cdot-n)} = (\th \star \bar\th)(n) = \de_{0,n} \eqdef
		\choice{
			1 \qifq n=0, \\
			0 \quad\text{otherwise.}
		}
	}
	where $\bar\th(x)=\overline{\th(-x)}$ and $\de_0$ is the discrete Dirac vector.
	 
By Tonelli's theorem and Plancherel's identity, $A\in L^1(\TT)$. Its Fourier coefficients are
	\begin{align*}
		\frac{1}{2\pi}\int_0^{2\pi}A(\om)e^{-\imath n\om}\,\d\om
		&=\frac{1}{2\pi}\int_\RR |\hat\th(\om)|^2e^{-\imath n\om}\,\d\om\\
		&=\dotp{\th}{\th(\cdot+n)}.
	\end{align*}
	Thus the translates are orthonormal exactly when the Fourier coefficients of $A$ are those of the constant function 1. Uniqueness of Fourier coefficients in $L^1(\TT)$ gives $A=1$ almost everywhere.
	 
	Since $A$ is $2\pi$-periodic, normalization by the bounded function $1/\sqrt{A(\om)}$ gives $\sum_{k} |\hat \phi(\om-2\pi k)|^2 = 1$. It remains to check equality of the closed spans. Approximate the periodic multipliers $1/\sqrt A$ and $\sqrt A$ by trigonometric polynomials in $L^2(\TT)$. The bounds $a\leq A\leq b$ show that the corresponding products with $\hat\th$ and $\hat\phi$ converge in $L^2(\RR)$. Inverting the Fourier transform expresses each generator as an $L^2$ limit of finite linear combinations of translates of the other.
\end{proof}

Spline interpolation provides a typical application. For example, cubic splines are generated by translates of a box-spline function $\th$, which is a compactly supported piecewise polynomial.

                                                   
\section{Multiresolution Detail Spaces}

The detail spaces are the orthogonal complements associated with the inclusions $V_j \subset V_{j-1}$. Since these subspaces are closed,
\eq{
	\foralls j, \quad W_j \text{ is such that } V_{j-1} = V_j \oplus^\bot W_j. 
}
The following diagram shows these nested approximation spaces and their orthogonal complements:
\begin{center}
	\image{wavelets}{.8}{emb